% TOP_PROBLEM: {"id":30007748,"problem_number":"LOCAL-30007748","title":"Mechanisms of neighborhood effects","category_id":21,"category":{"id":21,"name":"economics","display_name":"Economics","description":"Open economic theory statements, published research agendas, and proposed scoped specifications; question provenance and historical priority are qualified per record.","slug":"economics","order_index":21,"created_at":"2026-10-08T00:00:00Z"},"status":"research_question","external_url":null,"legacy_ids":[],"published":false,"statement":"Identify the separate local earnings contributions of school quality, peers, violence, and environmental conditions within a specified childhood neighborhood exposure.","economics":{"original_id":237,"field":"Urban, housing and spatial economics","kind":"identification","formalization_basis":"adapted_research_question","source_status":"explicit_open","priority_status":"earliest_found","priority_note":"Earliest explicit statement located in this audit; no exhaustive priority search or first-ever claim. The mathematical specification below is adapted, not quoted from the source.","earliest_verified_reference":{"authors":["Eric Chyn","Lawrence F. Katz"],"title":"Neighborhoods Matter: Assessing the Evidence for Place Effects","year":2021,"url":"https://lkatz.scholars.harvard.edu/sites/g/files/omnuum5961/files/lkatz/files/chyn_katz_jep_2021_all.pdf","doi":"10.1257/jep.35.4.197","locator":"Conclusion, printed pp. 216–217; introduction p. 200","evidence_summary":"The authors explicitly identify mechanism weights, general-equilibrium scaling of mobility programs, and effects of revitalization on preexisting residents as open issues."},"current_reference":{"authors":["Lawrence F. Katz"],"title":"Neighborhood Effects and Missing Markets for Opportunity","year":2026,"url":"https://www.nber.org/papers/w35419","doi":"10.3386/w35419","locator":"Abstract","evidence_summary":"The abstract discusses missing markets and mobility assistance; the older Chyn–Katz review explicitly establishes the scaling and mechanism research agendas."},"formalization_note":"The source explicitly identifies the underlying gap or agenda. Population, horizon, interventions, estimands, and auxiliary assumptions are proposed operational choices, not a canonical source theorem. Partial later answers are compatible with this scoped question; the citation alone does not prove absence of every subsequent solution. The 2021 Chyn–Katz review supplies an earlier verified agenda than the original catalogue’s 2023 or 2026 supporting citation. Mathematical scope was checked independently; added definedness, feasibility, or identification conditions belong to this proposed formalization."},"statusQualification":"Published question or research agenda verified at the cited locator. The mathematical specification is adapted for this catalogue and includes additional assumptions; source publication does not establish first-ever priority or certify this exact model remains unsolved. The source explicitly identifies the underlying gap or agenda. Population, horizon, interventions, estimands, and auxiliary assumptions are proposed operational choices, not a canonical source theorem. Partial later answers are compatible with this scoped question; the citation alone does not prove absence of every subsequent solution. The 2021 Chyn–Katz review supplies an earlier verified agenda than the original catalogue’s 2023 or 2026 supporting citation. Mathematical scope was checked independently; added definedness, feasibility, or identification conditions belong to this proposed formalization.","statusReviewedAt":"2026-10-08"}
% FIRST_POSTED: 2026-10-08
% ENTRY_KIND: statement_only
% !TeX program = lualatex
\documentclass[11pt]{article}
\usepackage{fontspec}
\usepackage[a4paper,margin=25mm]{geometry}
\usepackage{amsmath,amssymb,mathtools}
\usepackage{xurl}
\usepackage[hidelinks]{hyperref}
\newcommand{\sref}[2]{\hyperlink{source:#1:#2}{[S#2]}}
\setlength{\emergencystretch}{3em}
\begin{document}
% problemId:30007748
\section*{Mechanisms of neighborhood effects}\label{30007748}
\subsection{Definitions and mathematical statement}
Consider children aged six through twelve whose families receive randomized offers to move within one metropolitan area. Let neighborhood exposure be a vector \(x=(s,p,v,e)\), where \(s\) measures school quality, \(p\) peer composition, \(v\) exposure to violence, and \(e\) environmental quality using independently defined scales. Let \(Y_i(x)\) be adult earnings under a sustained six-year exposure profile. For baseline vector \(x_0\) and a feasible intervention changing only component \(j\), write \(x_0+\delta e_j\), where \(e_j\) is the corresponding unit vector. Define
\[\theta_j=\left.\frac{\partial E[Y_i(x_0+\delta e_j)]}{\partial\delta}\right|_{\delta=0}.\]
Require a feasible neighborhood of the component intervention and existence of the mean derivative; use one-sided contrasts at boundaries. Identify the relative contributions of these components and their interactions to the total move effect, using independent interventions or explicit valid instruments. A voucher lottery shifts several components together and ordinarily identifies a bundled effect rather than each \(\theta_j\). Specify exposure timing, household responses, and exclusion restrictions, and report bounds if component variation is insufficient. Evaluate whether results transport beyond the stated age range and metropolitan setting. The source explicitly asks about relative mechanism weights; the vector, local derivatives, and six-year intervention are proposed operational details.

\par\textbf{Formulation and scope.} The source explicitly identifies the underlying gap or agenda. Population, horizon, interventions, estimands, and auxiliary assumptions are proposed operational choices, not a canonical source theorem. Partial later answers are compatible with this scoped question; the citation alone does not prove absence of every subsequent solution. The 2021 Chyn–Katz review supplies an earlier verified agenda than the original catalogue’s 2023 or 2026 supporting citation. Mathematical scope was checked independently; added definedness, feasibility, or identification conditions belong to this proposed formalization.

\par\textbf{Open-question provenance.} An explicit question or research agenda was verified at the source locator. This does not by itself certify that every later version remains unresolved \sref{30007748}{1}.

\par\textbf{Historical priority.} Earliest explicit statement located in this audit; no exhaustive priority search or first-ever claim. The mathematical specification below is adapted, not quoted from the source.

\subsection{Short English statement}
Identify the separate local earnings contributions of school quality, peers, violence, and environmental conditions within a specified childhood neighborhood exposure.

\subsection{Sources}
\begin{itemize}\raggedright
\item[\textnormal{[S1]}]\hypertarget{source:30007748:1}{} Eric Chyn, Lawrence F. Katz. 2021. Neighborhoods Matter: Assessing the Evidence for Place Effects. Conclusion, printed pp. 216–217; introduction p. 200.
\url{https://lkatz.scholars.harvard.edu/sites/g/files/omnuum5961/files/lkatz/files/chyn_katz_jep_2021_all.pdf}.
DOI: 10.1257/jep.35.4.197.
{\small\textit{Earliest verified question/agenda reference; historical priority qualified below. The authors explicitly identify mechanism weights, general-equilibrium scaling of mobility programs, and effects of revitalization on preexisting residents as open issues.}}

\item[\textnormal{[S2]}]\hypertarget{source:30007748:2}{} Lawrence F. Katz. 2026. Neighborhood Effects and Missing Markets for Opportunity. Abstract.
\url{https://www.nber.org/papers/w35419}.
DOI: 10.3386/w35419.
{\small\textit{Supporting or current literature; see evidence qualification. The abstract discusses missing markets and mobility assistance; the older Chyn–Katz review explicitly establishes the scaling and mechanism research agendas.}}
\end{itemize}
\end{document}
